% Source: https://artofproblemsolving.com/community/c3h3525654_20point_connecting_lines_game_theory_auckland_mo_2024_p7
% Stable identity: aops-topic-3525654-q01
The first player joins two points which split the other eighteen points into two arcs of nine points each. Pair the points on one arc with those on the other in reversed cyclic order, as reflection across the first chord would do in a regular drawing. Any later legal chord has both endpoints in one open arc, since otherwise it would cross the first chord. After each opponent move, draw the corresponding chord on the other side. Reversal of cyclic order preserves the crossing and noncrossing relation between chords, so this response is always legal and uses unused endpoints. The two sides again have matching positions. Thus the first player can answer every move and makes the last move: the first player wins.
\begin{asy}
usepackage("amsmath");
import geometry;
size(260);
pair[] P;for(int k=0;k<20;++k){P.push(2*dir(18*k));dot(P[k]);}draw(circle((0,0),2));draw(P[0]--P[10],linewidth(1.2));draw(P[2]--P[6],blue);draw(P[18]--P[14],blue);label("$A$",P[0],E);label("$B$",P[10],W);label("$P$",P[2],NE);label("$Q$",P[6],NW);label("$P^{\prime}$",P[18],SE);label("$Q^{\prime}$",P[14],SW);
\end{asy}
